\documentclass[11pt,a4paper]{article}
\usepackage[utf8]{inputenc}
\usepackage[T1]{fontenc}
\usepackage[brazilian]{babel}
\usepackage[a4paper,margin=2.5cm]{geometry}
\title{Mapa de evolução -- Sprint 1}
\author{Matheus Moreli Durão Silva -- 12621CBS217}
\date{}
\begin{document}\maketitle
\section*{Uma entrega, construída por etapas}
Os arquivos desta pasta preservam primeiro o entregável obrigatório de cada card. Quando o exercício original não usa o tema do inventário, existe um \texttt{evolucao\_projeto.py} separado: ele aplica o mesmo conceito ao projeto sem substituir o código entregue. Assim, é possível estudar a atividade exigida e, ao lado dela, enxergar a evolução técnica até o sistema final.

\begin{enumerate}
\item \textbf{S1\_01 a S1\_03:} entendimento do problema, funções esperadas e critérios que mostram quando o sistema estará correto.
\item \textbf{S1\_04 a S1\_08:} exercícios básicos preservados; quando necessário, um complemento aplica \texttt{print}, variáveis, entrada, condições e repetição ao inventário.
\item \textbf{S1\_09 a S1\_13:} funções, listas, tuplas, sets, dicionários e \texttt{Enum}; a atividade original vem antes da relação com o projeto.
\item \textbf{S1\_14 e S1\_15:} salvamento em JSON e tratamento dos erros mais prováveis.
\item \textbf{S1\_16 a S1\_20:} histórico, colaboração, branches, revisão e conferência da versão final.
\end{enumerate}

\section*{Chegada ao projeto oficial}
Depois de praticar os blocos separados, a versão oficial reúne tudo no arquivo \texttt{app.py}, dentro da pasta \texttt{Projeto-Sprint-1} na Área de Trabalho. Ela acrescenta o fluxo completo: cadastro, consulta, atualização, remoção, vulnerabilidades, resumo e persistência. Os comentários presentes no arquivo explicam a razão de cada parte.

\section*{Como usar para estudar}
Abra uma pasta por vez, leia o PDF, execute o arquivo Python e altere apenas uma linha por teste. Em seguida, compare-o com o próximo card. Esse caminho demonstra tanto o checklist individual quanto a evolução técnica até o programa final.
\end{document}
